Courts are the one redistricting reform channel that does not require incumbent legislators to vote for change. When a court strikes a redistricting map, it can order a remedy---and that remedy can include an algorithmic redistricting requirement. This is the most powerful available bypass channel in states without ballot initiative mechanisms, and it is available in every state.

The judicial redistricting remedy has a long history. \textit{Reynolds v. Sims}, 377 U.S.\ 533 (1964), and its companion cases established that federal courts could order redistricting when states failed to comply with the Equal Protection Clause's one-person, one-vote requirement. Courts have since ordered redistricting remedies on racial grounds (\textit{Alabama Legislative Black Caucus v. Alabama}, 575 U.S.\ 254 (2015); \textit{Allen v. Milligan}, 599 U.S.\ 1 (2023)), on state constitutional grounds (\textit{League of Women Voters v. Pennsylvania}, 178 A.3d 737 (Pa. 2018); \textit{Harkenrider v. Hochul}, 38 N.Y.3d 494 (2022); \textit{Harper v. Hall}, 380 N.C.\ 317 (2022)), and on a combination of grounds.

Yet courts have not, to date, ordered algorithmic redistricting by name. Special masters in court-ordered redistricting typically apply the court's stated criteria through their own cartographic judgment---the same process that produces legislative gerrymanders, just with a neutral expert rather than an interested legislator. The special master's map is a product of individual judgment that opposing experts can challenge and that courts cannot independently verify.

The \texttt{bisect} platform offers courts a verifiable alternative. A map produced by the platform from public inputs can be reproduced byte-for-byte by any party running the platform with the recorded parameters. Courts can verify the algorithm's neutrality by inspecting its source code. Expert disputes about map quality can be resolved by reference to the manifest rather than by battles of cartographic experts.

This paper addresses legal practitioners and courts. Section~2 analyzes the special-master redistricting precedent and identifies the pathologies that algorithmic remedies could cure. Section~3 presents three model court orders, each appropriate for different procedural circumstances. Section~4 explains the SHA-256 audit chain and its significance for court reproducibility standards. Section~5 concludes.
